\begin{center}
    {\LARGE \textbf{2020分析化学}} \\[2ex]
    {\large 时量180分钟 \quad 满分90分}
\end{center}

\vspace{0.5cm}
\noindent\textbf{注意事项:}
\begin{enumerate}[label=\arabic*., parsep=0pt, itemsep=0pt]
    \item 答题前,考生先将自己的姓名、学号、考场号、座位号填写在答题卡上,将条形码准确粘贴在考生信息条形码粘贴区。
    \item 考试结束后,将本试卷与答题纸一并收回。
    \item 回答各个问题时,将答案写在答题纸上,写在本试卷上无效。
\end{enumerate}

\vspace{0.5cm}
\noindent\textbf{一、简答题(42 pts.)}

\begin{enumerate}[label=\arabic*., itemsep=1.5ex]
    \item 下列情况对分析结果的影响（填偏高、偏低或无影响）
    \begin{enumerate}[label=(\arabic*), noitemsep, topsep=0pt, partopsep=0pt]
        \item $\ce{NaOH}$溶液浓度标定后由于保存不妥吸收了$\ce{CO2}$，以此标准溶液测定草酸摩尔质量。
        \item $\ce{NaOH}$溶液浓度标定后由于保存不妥吸收了$\ce{CO2}$，若以此标准溶液测定$\ce{H3PO4}$浓度(甲基橙指示剂)。
    \end{enumerate}
    \item 某物质的水溶液$100\ \mathrm{mL}$，用$5$份$10\ \mathrm{mL}$萃取剂溶液连续萃取$5$次，总萃取率为$87\%$，则该物质在此萃取体系中的分配比是多少？
    \item 某生采用以下方法测定溶液中$\ce{Ca^{2+}}$的浓度，请指出其错误。先用自来水洗净$250\ \mathrm{mL}$锥形瓶，用$50\ \mathrm{mL}$容量瓶量取试液倒入锥形瓶中，然后加少许铬黑$\mathrm{T}$指示剂，用$\mathrm{EDTA}$标准溶液滴定。
    \item 求$[\ce{H^+}]=0.1\ \mathrm{mol/L}$的$\ce{H2SO4}$介质中$\ce{MnO4^- +5e^- +8H^+ = Mn^{2+} +4H2O}$半反应的条件电位。[已知$\varphi^\ominus_{\ce{MnO4^-/Mn^{2+}}}=1.51\ \mathrm{V}$]
    \item 什么是均匀沉淀法？其优点是什么？
    \item 写出两种标定盐酸溶液的基准物质。
    \item 某法测定铅锡合金中锡的质量分数，已知标准偏差为$0.08\%$，重复四次测定，质量分数$w(\ce{Sn})$的平均值为$62.32\%$，试计算出置信度$95\%$($u=1.96$)和置信度为$99\%$($u=2.58$)时平均值的置信区间，并结合计算结果说明置信区间与置信度有何关系？
\end{enumerate}

\vspace{0.5cm}
\noindent\textbf{二、计算题(48 pts.)}

\begin{enumerate}[label=\arabic*., itemsep=1.5ex]
    \item (10分) 用$0.10\ \mathrm{mol/L}\ \ce{HCl}$滴定$0.10\ \mathrm{mol/L}\ \ce{NaOH}$，而$\ce{NaOH}$试液中还含有$0.10\ \mathrm{mol/L}\ \ce{NaAc}$，计算: (1) 化学计量点及化学计量点前后$0.1\%$的$\mathrm{pH}$; (2) 滴定至$\mathrm{pH}=9.0$时的终点误差。[已知$\mathrm{p}K_a(\ce{HAc})=4.74$]
    \item (10分) 用$2\times10^{-2}\ \mathrm{mol/L}\ \mathrm{EDTA}$滴定同浓度$\ce{Zn^{2+}}$，$\ce{Al^{3+}}$混合液中的$\ce{Zn^{2+}}$，以磺基水杨酸($\ce{H2L}$)掩蔽$\ce{Al^{3+}}$。若终点时$\mathrm{pH}$为$5.5$，选二甲酚橙为指示剂。计算终点误差。[已知$\lg K(\ce{ZnY})=16.5$；$\lg K(\ce{AlY})=16.1$；$\mathrm{pH}=5.5$时,$\lg\alpha_{Y(\mathrm{H})}=5.5$,$\mathrm{pZn_{终}}$(二甲酚橙)$=5.7$；$\lg\alpha_{\ce{Al(L)}}=6.4$]
    \item (10分) 在$\mathrm{pH}=2.0$的含有$0.01\ \mathrm{mol/L}\ \mathrm{EDTA}$及$0.10\ \mathrm{mol/L}\ \ce{HF}$的溶液中，当加入$\ce{CaCl2}$使溶液中的$c(\ce{Ca^{2+}})=0.10\ \mathrm{mol/L}$时，问(1)$\mathrm{EDTA}$的存在对生成$\ce{CaF2}$沉淀有无影响? (2)能否产生$\ce{CaF2}$沉淀？(不考虑体积变化) [已知$\ce{HF}$的$K_a=6.6\times10^{-4}$,$K_{\mathrm{sp}}(\ce{CaF2})=2.7\times10^{-11}$,$\lg K(\ce{CaY})=10.69$,$\lg\alpha_{Y(\mathrm{H})}=13.51$]
    \item (10分) 以邻二氮菲(用$\ce{R}$表示)作显色剂用光度法测定$\ce{Fe}$,显色反应为:
    \[
    \ce{Fe^{2+} + 3R = FeR3} \quad \lg K(\ce{FeR3})=21.3
    \]
    \begin{enumerate}[label=(\arabic*), noitemsep, topsep=0pt, partopsep=0pt]
        \item 若过量显色剂的总浓度为$[R']=1.0\times10^{-4}\ \mathrm{mol/L}$，问在$\mathrm{pH}=3.0$时能否使$99.9\%$的$\ce{Fe^{2+}}$络合为$\ce{FeR3}$? [$\mathrm{pH}=3.0$时,$\lg\alpha_{R(\mathrm{H})}=1.9$]
        \item 称取$0.5000\ \mathrm{g}$试样，制成$100\ \mathrm{mL}$试液，从中移取$10.00\ \mathrm{mL}$显色后稀释为$50\ \mathrm{mL}$，用$1.0\ \mathrm{cm}$比色皿，在$510\ \mathrm{nm}$处测得透射比$T=60.0\%$。求试样中$\ce{Fe}$的质量分数。[已知摩尔吸光系数$\varepsilon_{510}=1.1\times10^4\ \mathrm{L/(mol\cdot cm)}$,$A_r(\ce{Fe})=55.85$]
    \end{enumerate}
    \item (8分) 氯化钡试样用重量法测定钡的质量分数$w(\ce{Ba})$，三次结果为（$\%$）：$56.10$、$56.15$、$56.09$；用络合滴定法测钡，三次结果为（$\%$）：$56.01$、$55.96$、$55.89$。问$95\%$显著水平时滴定法结果是否明显偏低？[单边$F_{0.05,2,2}=19.00$；双边$t_{0.05,4}=2.78$,$t_{0.10,4}=2.13$]
\end{enumerate}